\phantomsection
\chapter*{Hargai dan Jadilah Diri Sendiri}
\addcontentsline{toc}{section}{Hargai dan Jadilah Diri Sendiri -- Devant Fadhillah}
\penulis{Devant Fadhillah}


\awalcerpen{R}{usdi adalah siswa yang dari luar kelihatan biasa aja}. Datang ke sekolah, duduk, pulang. Tapi di dalam kepalanya sering ramai sendiri. Dia gampang ngerasa kurang. Bukan karena nilainya jelek terus, bukan juga karena nggak punya temen sama sekali. Cuma setiap liat orang lain, dia langsung bandingin.


Ada yang lebih rame dikerumunin. Ada yang lebih gampang diajak ngobrol. Ada yang tinggal muncul aja udah dapat perhatian. Rusdi liat itu berulang-ulang sampai dia ngerasa dirinya ketinggalan. Dia jarang nanya, ``Gue aslinya gimana?'' yang lebih sering muncul malah, ``Kenapa gue nggak kayak mereka?''


Suatu hari, salah satu temannya datang ke sekolah bawa barang baru. Bukan barang aneh, cuma baru. Tapi cukup bikin orang-orang deketin. Ada yang nanya, ada yang pegang-pegang, ada yang cuma berdiri deket biar keliatan ikut. Rusdi liat dari kejauhan. Dia nggak ikut nyamperin. Cuma berdiri, nonton, terus ngerasa iri naik pelan-pelan.


Di kepala dia langsung ada pikiran yang sederhana. Kalau dia punya barang itu, mungkin orang-orang juga bakal nengok. Mungkin dia nggak perlu berdiri jauh-jauh. Mungkin dia juga bisa jadi yang dikerumunin.


Sepulang sekolah Rusdi langsung pulang, langsung minta. ``Beliin aku ini dong. Orang lain pada punya.'' Nadanya agak ngegas. Bukan minta pelan. Lebih ke nuntut. Orang tuanya udah sering denger permintaan kayak gitu, \textit{plus} alasan yang sama orang lain punya, jadi dia juga harus punya. Kali ini mereka nggak nurutin.

``Rus, kamu nggak harus punya semua yang orang lain punya,'' kata ibunya.


Rusdi kesal. Kesal yang bikin dia nggak mau denger lanjutan. Dia masuk kamar, pintu dibanting, terus duduk sendirian. Di situ dia mulai mikir lagi. Bukan soal barangnya doang. Lebih ke cara biar orang-orang mulai perhatiin dia. ``Gimana caranya supaya gue bisa kayak mereka?''


Malam itu Rusdi nonton film, awalnya cuma biar kamar nggak sepi. Dari film itu dia nangkep satu hal. Ada tokoh yang berpura-pura jadi orang lain supaya dapat perhatian dan kepercayaan. Tokoh itu nggak langsung kelihatan hebat dari awal. Dia nyalin. Dia pura-pura. Terus orang-orang mulai deket.


Rusdi dapat ide. Awalnya masih ngerasa itu cuma pikiran iseng. ``Kalau gue jadi orang lain gimana, ya?'' Kayak candaan di kepala. Tapi keesokan harinya ide itu dicoba.


Dia mulai ngecek orang yang menurutnya lebih keren. Cara ngomongnya ditiru. Cara berpakaiannya dideketin. Cara bergaulnya ikut-ikutan. Kalau orang itu nyapa dengan gaya tertentu, Rusdi coba gaya yang sama. Kalau orang itu kelihatan santai di keramaian, Rusdi maksa dirinya kelihatan santai juga.


Beberapa orang mulai nengok. Ada yang nyamperin. Ada yang ngajak ngobrol. Rusdi seneng. Akhirnya ada yang liat. Akhirnya ada yang ngeanggap dia ada. Perhatian yang dari dulu dia cari, sekarang keliatan masuk.


Masalahnya, rasa itu bikin dia nggak mau berhenti. Lama-lama meniru nggak cukup. Dia mulai nambah-nambahin cerita soal dirinya. Biar kelihatan lebih hebat. Dia mulai bilang punya banyak uang. Bilang bisa beli apa aja. Padahal kenyataannya nggak gitu.


Bohong kecil itu butuh penutup. Rusdi takut ketahuan. Dari situ dia nekat. Dia ingat ide yang muncul setelah nonton film, lalu bikin uang palsu dan nyoba pake. Pas uang itu sempat diterima, dia ngerasa yakin. Kayaknya caranya berhasil. Kayaknya dia bisa terus jalanin hidup kayak orang yang dia tiru.


Tapi rasa yakin itu nggak bertahan lama. Suatu hari, pas Rusdi lagi pake uang palsu itu, perbuatannya ketahuan. Kabarnya sampe ke orang tua. Setelah itu, kebohongan-kebohongan lain ikut kebuka. Bukan cuma soal uang. Soal siapa dia selama ini juga.


Waktu ditanya kenapa dia ngelakuin semua itu, Rusdi diam dulu. Yang keluar akhirnya jujur, pelan. ``Aku cuma pengen jadi orang lain yang selalu diperhatikan. Aku pengen orang-orang liat aku juga.''


Ayahnya tarik napas. ``Kamu pikir mereka memperhatikan kamu karena kamu jadi orang lain?'' Rusdi nggak jawab. Ibunya nyusul, ``Kalau kamu terus meniru orang lain, kapan kamu belajar menghargai diri kamu sendiri?''


Pertanyaan itu bikin Rusdi terdiam lama. Bukan karena dia tiba-tiba jadi bijak. Lebih karena baru kerasa, selama ini dia terlalu sibuk ngejar hidup orang lain. Dia kira kalau berhasil meniru seseorang, dia bakal dihargai. Padahal makin jauh dia meniru, makin banyak yang harus ditutupi. Makin banyak juga bagian dirinya yang ditinggal.


Rusdi akhirnya terima akibatnya. Dia minta maaf ke orang tuanya. Minta maaf juga ke orang-orang yang rugi karena tindakannya. Nggak ada yang langsung bikin semua terasa ringan. Cuma paling nggak dia nggak lari lagi dari yang udah terjadi.


Beberapa hari kemudian Rusdi balik ke sekolah. Kali ini dia nggak nyiapin versi lain dari dirinya. Dia pake barang yang emang dia punya. Ngomong dengan cara dia sendiri. Nggak dipaksa biar keliatan hebat. Perhatiannya nggak sebanyak waktu dia lagi pura-pura. Tapi itu lebih gampang ditahan daripada takut ketahuan.


Iri masih muncul. Pas liat orang lain punya sesuatu yang dia nggak punya, kepala dia masih kadang bandingin. Cuma sekarang dia coba nahan. Dia inget akibatnya. Inget pertanyaan ibunya. Inget betapa capeknya jaga kebohongan.


Di perjalanan pulang, Rusdi senyum kecil. Bukan senyum puas banget. Lebih ke senyum yang capek tapi lega. Dia bilang ke dirinya sendiri, ``Gue nggak harus jadi apa pun atau siapa pun cuma buat bisa dihargai.''


Sejak itu Rusdi mulai belajar pelan-pelan. Jadi diri sendiri emang nggak gampang. Kadang masih ada rasa pengen diperhatikan. Kadang masih ada godaan buat meniru orang. Tapi setidaknya sekarang dia tau, nggak ada gunanya ilangin diri sendiri cuma karena pengen dapat perhatian dari orang lain.

\jedahias
